\section{Limitations}
In our internal tests, we have identified some limitations in current Kimi K2 models. When dealing with hard reasoning tasks or unclear tool definition, the model may generate excessive tokens, sometimes leading to truncated outputs or incomplete tool calls. Additionally, performance may decline on certain tasks if tool use is unnecessarily enabled. When building complete software projects, the success rate of one-shot prompting is not as good as using K2 under an agentic coding framework. We are working to address these issues in future releases and looking forward to more feedbacks.

\section{Conclusions}

We introduced Kimi K2, a 1T-parameter open-weight MoE model built for agentic intelligence. Leveraging the token-efficient MuonClip optimizer and a 15.5T-token high-quality dataset, Kimi K2 achieves stable, scalable pre-training. Post-training combines large-scale synthetic tool-use data with a unified RL framework using both verifiable rewards and self-critic feedbacks. Kimi K2 sets new state-of-the-art on agentic and reasoning benchmarks, establishing itself as the most capable open-weight LLM to date.

